\documentclass[11pt]{article}
\usepackage[a4paper,margin=2.4cm]{geometry}
\usepackage[T1]{fontenc}
\usepackage{lmodern}
\usepackage{microtype}
\usepackage{graphicx}
\usepackage{booktabs}
\usepackage{siunitx}
\usepackage{xcolor}
\usepackage{amssymb}
\usepackage{enumitem}
\usepackage{caption}
\usepackage[colorlinks=true,linkcolor=accent,urlcolor=accent]{hyperref}

\definecolor{accent}{HTML}{2B5F6B}
\definecolor{good}{HTML}{2E7D4F}
\definecolor{warn}{HTML}{A9791F}
\definecolor{poor}{HTML}{A63A3A}
\definecolor{faint}{HTML}{F0F3F5}

\captionsetup{font=small,labelfont=bf,labelsep=period}
\setlist{leftmargin=1.4em,itemsep=2pt,topsep=3pt}
\newcommand{\g}[1]{\textcolor{good}{#1}}
\newcommand{\w}[1]{\textcolor{warn}{#1}}
\newcommand{\bd}[1]{\textcolor{poor}{#1}}
\sisetup{table-format=1.3,detect-weight=true}

% callout box
\usepackage[most]{tcolorbox}
\newtcolorbox{callout}[1][accent]{colback=faint,colframe=#1,boxrule=0pt,
  leftrule=3pt,arc=2pt,left=8pt,right=8pt,top=6pt,bottom=6pt}

\title{\vspace{-1.4cm}\bfseries Recognising individual bats by face:\\[2pt]
\large what the multi-seed results actually show}
\author{\normalsize Master's thesis --- results interpretation}
\date{\normalsize Open-set individual recognition \textbullet\ 2 species \textbullet\ 3 model families
\textbullet\ 3 backgrounds \textbullet\ 5 seeds/cell (270 runs)}

\begin{document}
\maketitle
\vspace{-0.8em}

\begin{callout}
\textbf{In one paragraph.} Across a re-run of the full experiment matrix, individual bat faces are
recognisable well above chance in the favourable conditions (best cells ROC-AUC $\approx$
0.85--0.89; permutation $p$ from $<0.001$ to $0.007$ on the thresholded operating point), and the signal is genuinely in the
\emph{foreground} --- performance survives removing the background entirely. Two limits are equally
clear: the models are not background-invariant (randomised backgrounds drop them to near chance),
and higher input resolution does not help (224$\to$320\,px, mean $\Delta$ROC-AUC $=-0.033$).
Metric-learning embeddings (ArcFace/AdaFace) are more robust than the contrastive Siamese network,
especially where training identities are few. A control experiment adds a decisive qualifier: much
of the apparent recognition --- nearly \textbf{all} of it for Rousettus --- is reproducible from the
crop \textbf{silhouette} alone, so the signal is as much head-shape as facial texture. The dominant
caveat throughout is the tiny open-set test (two held-out identities per species), so every claim is
directional, not definitive.
\end{callout}

\section{What the experiment measured}
The task is \textbf{open-set} individual recognition: the identities used to test the model were
never seen in training. Each species' data is split by identity, and the held-out \emph{test} set is
small (Table~\ref{tab:splits}) --- the single most important fact for interpreting everything below.

\begin{table}[h]\centering
\caption{Identity-disjoint splits.}\label{tab:splits}
\begin{tabular}{lccccc}
\toprule
Species & Train id. & Val id. & Test id. & Train img & Test img \\
\midrule
Mauritius & 12 & 2 & \bd{2} & 373 & 77 \\
Rousettus & 8 & 2 & \bd{2} & 619 & 137 \\
\bottomrule
\end{tabular}
\end{table}

Three model families were compared: a \textbf{Siamese} 4-conv contrastive network (learns a
similarity from pairs) and two margin-based metric-learning heads on a ResNet-50 backbone,
\textbf{ArcFace} and \textbf{AdaFace} (learn an identity-structured embedding). Each was trained on
three \textbf{background variants} of the same face crops --- \emph{original} (natural background
kept), \emph{green} (face segmented onto a uniform field, i.e.\ background removed), and \emph{random}
(face composited onto varied, identity-unrelated backgrounds) --- at five seeds each. Verification
quality is ROC-AUC; identification is top-1/mAP; a \textbf{permutation test} checks each result
against a shuffled-label null.

\paragraph{How to read the numbers.} With only two test identities, a single run is noisy: one
pairing can swing ROC-AUC by $>0.1$. Every cell is therefore a \textbf{mean $\pm$ standard deviation
over five seeds}, and chance is ROC-AUC $=0.5$. Differences smaller than roughly one standard
deviation are ties. The permutation test asks the stricter question --- could this arise from a model
that learned nothing? --- and for the strong cells the answer is no ($p$ from $<0.001$ to $0.007$).
These p-values are computed on the thresholded operating point (F1 and accuracy/precision/recall),
not on ROC-AUC itself, over $B=1000$ label shuffles. The smallest value $B$ can resolve is
$1/(B+1)\approx0.001$, so cells reported at that value are \emph{censored} at the floor rather than
estimated: the truth lies at or below it.

\section{The headline results}
\begin{table}[h]\centering
\caption{Test ROC-AUC (mean $\pm$ std, 5 seeds), edge 112\,px.
\g{Green} $\ge 0.75$, \w{amber} 0.60--0.75, \bd{red} $<0.60$ (approaching chance).}\label{tab:matrix}
\begin{tabular}{ll cc}
\toprule
Model & Background & Mauritius & Rousettus \\
\midrule
Siamese & original & \g{$0.880\pm.059$} & \w{$0.614\pm.106$} \\
Siamese & green    & \g{$0.798\pm.105$} & \bd{$0.421\pm.051$} \\
Siamese & random   & \w{$0.700\pm.050$} & \bd{$0.454\pm.035$} \\
\addlinespace
ArcFace & original & \g{$0.885\pm.068$} & \w{$0.740\pm.119$} \\
ArcFace & green    & \g{$0.777\pm.122$} & \g{$0.774\pm.085$} \\
ArcFace & random   & \w{$0.611\pm.073$} & \bd{$0.508\pm.135$} \\
\addlinespace
AdaFace & original & \g{$0.845\pm.071$} & \w{$0.739\pm.056$} \\
AdaFace & green    & \g{$0.809\pm.114$} & \g{$0.757\pm.127$} \\
AdaFace & random   & \bd{$0.592\pm.061$} & \bd{$0.551\pm.102$} \\
\bottomrule
\end{tabular}
\end{table}

\begin{figure}[h]\centering
\includegraphics[width=0.95\linewidth]{figures/multiseed_matrix.png}
\caption{Test ROC-AUC (mean $\pm$ std) by background, per species; dashed line is chance (0.5).
Two patterns recur: \emph{random} backgrounds are lowest everywhere, and Rousettus is both lower
and more variable than Mauritius.}
\end{figure}

\section{Finding 1 --- the recognition is real}
The concern with a two-identity test is that a good score is a fluke. The permutation test re-scores
each model against 1000 label shuffles; for the strong cells (Mauritius ArcFace/AdaFace on
original background) the observed operating-point score beats \emph{every} shuffle, giving the
smallest p-value $B=1000$ can express ($p<0.001$); the weaker of the strong cells reach $p=0.007$.
\begin{callout}[good]
\textbf{Defensible claim.} In the favourable conditions the models discriminate individuals at a rate
shuffled-label chance does not reproduce --- the learned representation carries genuine identity
information.
\end{callout}

\section{Finding 2 --- foreground, not background (but not background-invariant)}
If the high scores were the model memorising the enclosure rather than the bat, removing the
background should destroy performance. It does not:
\begin{itemize}
\item \textbf{Background removed (green):} embedding models still recognise well --- ArcFace
0.78\,/\,0.77 and AdaFace 0.81\,/\,0.76 (Mauritius\,/\,Rousettus), on par with the natural background.
\item \textbf{Background randomised (random):} the same models collapse toward chance (0.51--0.61).
\end{itemize}
Green working while random fails is the tell: a clean uniform field lets the network read the
foreground; arbitrary cluttered backgrounds act as distractors it was never taught to suppress. So
the model has learned a foreground representation that is \emph{not robust} to background variation.
\emph{Which part} of the foreground it uses --- the face, or the crop shape --- is the subject of
Section~\ref{sec:validity}.

\begin{figure}[h]\centering
\includegraphics[width=0.49\linewidth]{figures/proj_mauritius_umap.png}\hfill
\includegraphics[width=0.49\linewidth]{figures/proj_rousettus_umap.png}
\caption{UMAP of ArcFace test embeddings, coloured by identity (left Mauritius, right Rousettus,
original background). Tight, separated colour clusters mean same-individual images sit together and
different individuals are pushed apart --- the geometric basis of the ROC-AUC numbers.}
\end{figure}

\section{Finding 3 --- embeddings beat Siamese, most where identities are scarce}
The families diverge sharply on Rousettus (fewer training identities, 8 vs 12): the Siamese network
barely functions (0.61 original, \textbf{0.42 green --- below chance}), while ArcFace/AdaFace hold
0.74--0.77. The Mauritius$\to$Rousettus drop is therefore largest for Siamese (green: 0.80$\to$0.42).
The contrastive objective learns a metric from pairs and needs a diversity of identities to
generalise; with only eight it underfits. Margin-based heads extract more per identity and transfer
better to the open-set test.

\begin{figure}[h]\centering
\includegraphics[width=0.49\linewidth]{figures/siam_mauritius_umap.png}\hfill
\includegraphics[width=0.49\linewidth]{figures/siam_rousettus_umap.png}
\caption{Siamese embeddings, UMAP (left Mauritius, right Rousettus, original) --- the family with the
largest species gap. Mauritius shows structure; Rousettus is diffuse, the geometric counterpart of
its near-chance ROC-AUC. (This is why the resolution study excludes Siamese: its 4-conv net is fixed
at 105\,px.)}
\end{figure}

\section{Finding 4 --- higher resolution does not help}
Retraining the embedding models on genuinely higher-resolution crops at 224 and 320\,px did not
improve recognition: mean change 224$\to$320 was $-0.033$ ROC-AUC (green $-0.048$, original $-0.036$,
random $-0.015$). GradCAM localises the eyes \emph{more} sharply at higher resolution
(Fig.~\ref{fig:gradcam}), so the model is not starved of detail --- it already has what it needs at
112\,px.
\begin{figure}[h]\centering
\includegraphics[width=0.9\linewidth]{figures/gradcam_resolution_mauritius_class_logit.png}
\caption{GradCAM (class-logit target) at 112/224/320\,px, Mauritius. Attention concentrates on the
facial region and sharpens with resolution --- yet accuracy is unchanged.}\label{fig:gradcam}
\end{figure}

\section{Internal validity: is it the face, or the crop shape?}\label{sec:validity}
\begin{callout}
\textbf{Two terms, used throughout this section.}\\
\textbf{Shape (silhouette)} --- the outline of the segmented head crop only: head size, ear notches,
jaw contour. The black-vs-white boundary in the ``Silhouette'' panels (Figs.~\ref{fig:montage-mau}--\ref{fig:montage-rou}),
with all interior detail erased.\\[3pt]
\textbf{Texture} --- shorthand for everything \emph{inside} that outline: the eyes (position, spacing,
iris), the snout and muzzle patterning, fur colour and markings, shading --- the internal facial
appearance a person would use to tell two bats apart. (``Texture'' is the standard computer-vision
term; it means internal appearance here, not just fine fur detail.)
\end{callout}

Finding~2 showed recognition survives removing the background (green). But green keeps the tight
segmentation \textbf{silhouette} --- the head/face outline --- which is itself identity-correlated
(the ``Silhouette'' panels: head size, ear notches and jaw contour visibly differ between
individuals). So ``green works'' might reflect the model reading the \emph{shape} of the crop, not the
\emph{texture} of the face. Two controls test this directly.

\begin{figure}[h]\centering
\includegraphics[width=\linewidth]{figures/montage_mauritius.png}
\caption{Mauritius --- the same test crop as four inputs: natural, background-removed (green),
\textbf{silhouette} (shape only, zero texture), and \textbf{aperture} (texture only, outline removed).
The silhouette carries no facial detail, yet is trainable and discriminative.}\label{fig:montage-mau}
\end{figure}
\begin{figure}[h]\centering
\includegraphics[width=\linewidth]{figures/montage_rousettus.png}
\caption{Rousettus --- same four inputs. Note how distinctive the bare outline is even with the face
erased.}\label{fig:montage-rou}
\end{figure}

\subsection*{Control 1 --- train on shape alone (silhouette)}
Retrain the full matrix on the pure silhouette, with identical identities and splits, and compare to
green (shape\,+\,texture):
\begin{table}[h]\centering
\caption{Silhouette (shape only) vs green, test ROC-AUC (mean $\pm$ std, 5 seeds).
\bd{Red $\Delta$} = silhouette matches/beats green (shape dominates); \g{green $\Delta$} = texture
adds value.}\label{tab:silh}
\begin{tabular}{ll ccc}
\toprule
Species & Model & Silhouette & Green & $\Delta$ (silh$-$green) \\
\midrule
Mauritius & Siamese & $0.697\pm.023$ & 0.798 & \g{$-0.101$} \\
Mauritius & ArcFace & $0.649\pm.086$ & 0.777 & \g{$-0.128$} \\
Mauritius & AdaFace & $0.611\pm.130$ & 0.809 & \g{$-0.198$} \\
\addlinespace
Rousettus & Siamese & $0.698\pm.032$ & 0.421 & \bd{$+0.277$} \\
Rousettus & ArcFace & $0.823\pm.038$ & 0.774 & \bd{$+0.049$} \\
Rousettus & AdaFace & $0.820\pm.075$ & 0.757 & \bd{$+0.063$} \\
\bottomrule
\end{tabular}
\end{table}

Shape alone is \textbf{well above chance everywhere} (0.61--0.82). For Mauritius the silhouette runs
0.10--0.20 \emph{below} green, so internal texture genuinely adds value. For Rousettus the silhouette
\textbf{matches or beats} green --- Siamese by $+0.28$ --- so the bare outline is as discriminative as
the textured face, or more.

\subsection*{Control 2 --- occlude one cue on the real-face model (5 seeds)}
\begin{table}[h]\centering
\caption{Occluding one cue in the test set of a model trained on real faces (ArcFace/original, mean
over 5 seeds). Bigger drop $=$ stronger reliance on the removed cue.}\label{tab:occl}
\begin{tabular}{ll cc}
\toprule
Species & Test-set ablation & ROC-AUC & $\Delta$ vs clean \\
\midrule
Rousettus & clean & $0.793\pm.135$ & --- \\
Rousettus & shape removed (texture only) & $0.672\pm.117$ & \bd{$-0.122$} \\
Rousettus & texture removed (shape only) & $0.774\pm.141$ & $-0.019$ \\
\addlinespace
Mauritius & clean & $0.808\pm.185$ & --- \\
Mauritius & shape removed (texture only) & $0.884\pm.047$ & \textcolor{gray}{$+0.076$ (noise)} \\
Mauritius & texture removed (shape only) & $0.815\pm.101$ & \textcolor{gray}{$+0.007$ (noise)} \\
\bottomrule
\end{tabular}
\end{table}

For Rousettus, removing the \emph{shape} costs $6\times$ more than removing the \emph{texture}
($-0.122$ vs $-0.019$) --- it relies on shape, corroborating Control~1. For Mauritius the deltas are
swamped by the two-identity test's $\pm0.20$ noise and are inconclusive (occlusion even nominally
``improves'' it --- a sign of pure variance).

\subsection*{Why this is not negligible}
\begin{itemize}
\item \textbf{It is most of the signal, not a small correction.} For Rousettus a binary silhouette
with zero facial detail reaches ROC-AUC 0.82 (ArcFace) --- \emph{higher} than the textured face
(0.77). Nearly all of the reported Rousettus recognition is reproducible from the outline alone.
\item \textbf{Two independent controls agree.} For Rousettus, silhouette $\ge$ green (Control~1)
\emph{and} shape-removal hurts most (Control~2).
\item \textbf{It changes what the result means.} A ``face recognition'' score that survives deleting
the face --- keeping only the crop outline --- is measuring head/crop \emph{geometry}, not facial
identity. The silhouette is the segmentation mask, whose contour (head size, ear shape, jaw line) is
stable within an individual and variable across individuals: a per-individual shortcut a CNN will
exploit.
\item \textbf{It is species-specific.} For Mauritius texture \emph{does} contribute (silhouette drops
below green), so the face claim is partly supported --- though shape still accounts for a large share.
For Rousettus a distinct facial-texture contribution is not established by these data.
\end{itemize}
\begin{callout}[warn]
\textbf{State it carefully.} Report the silhouette baseline next to every headline number. Honest
phrasing: Mauritius --- ``individuals are recognised from face crops, with contributions from both
facial texture and head-crop shape''; Rousettus --- ``individuals are separable largely by head-crop
\textbf{shape}; a distinct facial-texture contribution is not demonstrated.'' Isolating texture
cleanly would require shape-normalised inputs (a fixed aperture) at training time, not just at test
time.
\end{callout}

\section{What the thesis learns}
\begin{itemize}
\item \textbf{Feasibility.} Automated individual identification of these fruit-bat species from face
crops is possible above chance --- a proof of concept, not yet a deployable tool.
\item \textbf{Architecture.} Prefer margin-based embedding learning (ArcFace/AdaFace) over contrastive
Siamese, especially with few identities.
\item \textbf{The real bottleneck is data.} Neither model capacity nor input resolution is limiting;
the scarce resource is the \emph{number of individuals}.
\item \textbf{Background must be controlled}, and \textbf{shape is a confound, not a footnote} --- a
silhouette baseline recognises individuals on its own, so the ``face'' framing must be qualified and
the silhouette number reported beside every accuracy.
\end{itemize}

\section{Limitations to state up front}
\begin{itemize}
\item \textbf{Two-identity open-set test} --- wide confidence intervals; conclusions are about
patterns across the matrix, not any single cell.
\item \textbf{Background robustness gap} --- the random-background collapse means the model would not
transfer to unfamiliar backgrounds without further work.
\item \textbf{Shape confound} --- a large share of performance (nearly all for Rousettus) is
reproducible from the crop silhouette; isolating texture needs shape-normalised training inputs.
\item \textbf{Generalisation scope} --- splits are identity-disjoint but within the same capture
setting; cross-session generalisation is not separately demonstrated.
\end{itemize}

\begin{callout}[good]
\textbf{For the abstract.} ``On an open-set test of held-out individuals, margin-based face embeddings
recognise individual fruit bats above chance (best ROC-AUC $\approx$ 0.85--0.89; permutation
$p<0.01$), with the signal shown to be foreground- rather than background-driven --- though a
silhouette control indicates much of it (nearly all for Rousettus) reflects head-crop \emph{shape}
rather than facial texture; performance is limited by the number of available identities rather than
by model capacity or image resolution, and is not yet robust to background change.''
\end{callout}

\end{document}
